\documentclass[runningheads]{llncs}
\usepackage[T1]{fontenc}
\usepackage{graphicx}
\usepackage{booktabs}
\usepackage{hyperref}
\usepackage{amsmath}

\begin{document}

\title{Supplier Lead-Time Risk Prediction for Retail Procurement}

\author{Nguyen Hoai Khanh\inst{1} \and
Ho Lam Bao Dang\inst{1} \and
Duong Gia Bao\inst{1}}

\authorrunning{Khanh et al.}

\institute{FPT University HCMC, Ho Chi Minh City, Vietnam\\
\email{\{K@itsdre.me, baodang10506@gmail.com, ilyrabao@gmail.com\}}}

\maketitle

\begin{abstract}
Supplier lead-time variability poses significant stockout and procurement risks in
retail operations. This study builds a predictive classification framework using
supply chain logistics data, benchmarking four classifiers (Logistic Regression,
Decision Tree, Random Forest, XGBoost) with SHAP explainability to identify
high-risk suppliers. Primary evaluation uses PR-AUC and recall on the delayed
class to account for class imbalance (16\% delayed). XGBoost achieves the highest
PR-AUC under stratified 5-fold cross-validation. A Streamlit web application
enables procurement teams to score individual shipments in real time.

\keywords{Supplier Risk Prediction \and Delay Classification \and Supply Chain Analytics \and XGBoost \and SHAP}
\end{abstract}


\section{Introduction}

Supplier lead-time delays cause stockouts, margin erosion, and poor customer
satisfaction in retail procurement. Early identification of high-risk shipments
allows procurement teams to pre-emptively switch carriers or expedite orders.

This project addresses three research questions:
\begin{enumerate}
  \item Which suppliers and routes carry the highest delay risk?
  \item Which machine learning model best predicts shipment delays?
  \item How can predictions be explained and acted on by non-technical users?
\end{enumerate}


\section{Dataset}

The analysis uses three publicly available CSV files from Kaggle, merged into a
single modelling dataset of \textbf{704 rows and 23 features}.

\begin{table}[h]
\caption{Data sources}\label{tab:data}
\begin{tabular}{lrl}
\toprule
File & Rows & Contents \\
\midrule
\texttt{customer.csv}              & 750 & Supplier attributes, lead-time days, market segment \\
\texttt{shipment.csv}              & 728 & Delivery status (target), customs time, freight cost \\
\texttt{logistics\_performance.csv} & 100 & Carrier delay averages, warehouse utilisation by region \\
\bottomrule
\end{tabular}
\end{table}

\texttt{shipment.csv} serves as the fact table. \texttt{customer.csv} is joined
on \texttt{supplier\_id}; \texttt{logistics\_performance.csv} is joined on
\texttt{region} (derived from destination country). Inner joins reduce the
combined source rows to the final 704-row modelling set.

The target is whether a shipment was \textbf{Delayed} (16\%) or \textbf{On-Time}
(84\%). Class imbalance is addressed through \texttt{class\_weight='balanced'} in
model training and threshold tuning during evaluation.


\section{Methods}

\subsection{Feature Engineering}

Categorical columns (carrier, market segment, region) are label-encoded.
Numeric features are standardised with \texttt{StandardScaler}. The dataset is
split 80/20 with stratification on the target.

\subsection{Models}

Four classifiers are benchmarked:
\begin{itemize}
  \item Logistic Regression — baseline; coefficients interpreted as odds ratios
  \item Decision Tree
  \item Random Forest
  \item XGBoost — primary model deployed in the Streamlit application
\end{itemize}

\subsection{Evaluation}

\textbf{Primary:} PR-AUC and recall on the Delayed class (imbalance-aware).\\
\textbf{Secondary:} ROC-AUC, F1-score, confusion matrix.\\
\textbf{Validation:} Stratified 5-fold cross-validation.

\subsection{Explainability}

SHAP TreeExplainer is applied to the XGBoost model to produce global feature
importance rankings and per-shipment waterfall plots surfaced in the web
application.


\section{SQL Analysis}

SQLite queries compute supplier-level risk patterns:
\begin{itemize}
  \item Average lead time and delay frequency per supplier
  \item Volume--delay correlation
  \item Monthly delay trends and year-over-year growth rate
  \item Carrier penetration index (delay share vs volume share)
\end{itemize}


\section{Results}

% TODO: fill in after running 05_modeling.py
\begin{table}[h]
\caption{Model comparison}\label{tab:results}
\begin{tabular}{lcccc}
\toprule
Model & PR-AUC & Recall (Delayed) & ROC-AUC & F1 \\
\midrule
Logistic Regression & -- & -- & -- & -- \\
Decision Tree       & -- & -- & -- & -- \\
Random Forest       & -- & -- & -- & -- \\
XGBoost             & -- & -- & -- & -- \\
\bottomrule
\end{tabular}
\end{table}


\section{Visualisation}

Plotly and Folium are used to produce:
\begin{itemize}
  \item Risk heatmap surfacing highest-risk supplier--route pairs
  \item Boxplots of lead-time spread per supplier
  \item Lead-time histograms and monthly trend lines
\end{itemize}

A Power BI supplier scorecard provides drill-down by supplier, region, and
carrier with automated risk alerts to support procurement decisions.


\section{Web Application}

The Streamlit application (\texttt{app/app.py}) accepts shipment input parameters
and returns a real-time delay probability from the trained XGBoost model alongside
a SHAP waterfall chart explaining the individual prediction.


\section{Conclusion}

% TODO: fill in after experiments complete

\bibliographystyle{splncs04}
\bibliography{references}

\end{document}
